\documentclass[tikz,border=10pt]{standalone}
\usepackage[no-math]{fontspec}
\usepackage[brazil]{babel}
\usepackage{amsmath,amssymb}
\IfFileExists{C:/Windows/Fonts/Lato_Regular.ttf}{%
  \setmainfont{Lato}[
    Path = C:/Windows/Fonts/,
    UprightFont = Lato_Regular.ttf,
    BoldFont = Lato_Bold.ttf,
    ItalicFont = Lato_Italic.ttf,
    BoldItalicFont = Lato_BoldItalic.ttf
  ]
  \setsansfont{Lato}[
    Path = C:/Windows/Fonts/,
    UprightFont = Lato_Regular.ttf,
    BoldFont = Lato_Bold.ttf,
    ItalicFont = Lato_Italic.ttf,
    BoldItalicFont = Lato_BoldItalic.ttf
  ]
}{%
  \RequirePackage[default]{lato}
}
\renewcommand{\familydefault}{\sfdefault}

\usepackage{xcolor}
\usepackage{tikz}
\usepackage{pgfplots}
\pgfplotsset{compat=1.18}
\usetikzlibrary{arrows.meta,positioning,shapes.geometric,fit,backgrounds,calc,decorations.pathreplacing,patterns}

% Paleta institucional PLANTERR / UEFS (alinhada ao estilo planterr-slides.css / 09a)
\definecolor{etNavy}{HTML}{1F4268}
\definecolor{etNavyDark}{HTML}{152E4A}
\definecolor{etGold}{HTML}{B38A3D}
\definecolor{etGoldDark}{HTML}{8C6927}
\definecolor{etTeal}{HTML}{0E7C7B}
\definecolor{etForest}{HTML}{2D6A4F}
\definecolor{etAlert}{HTML}{9E2A2B}
\definecolor{etInk}{HTML}{26313D}
\definecolor{etMuted}{HTML}{5B6573}
\definecolor{etBorder}{HTML}{D5DEE8}

\definecolor{bgNavy}{HTML}{EDF2F7}
\definecolor{bgGold}{HTML}{F8F3E9}
\definecolor{bgTeal}{HTML}{EBF6F5}
\definecolor{bgForest}{HTML}{EEF5F0}
\definecolor{bgAlert}{HTML}{FDF2F2}
\definecolor{bgCanvas}{HTML}{FFFFFF}

\tikzset{
  >=Stealth,
  etBox/.style={
    draw=#1, fill=white, rounded corners=5pt, line width=1.15pt,
    inner sep=7pt, align=center, font=\footnotesize\sffamily, text=etInk
  },
  etTint/.style={
    draw=#1, fill=#1!8!white, rounded corners=5pt, line width=1.15pt,
    inner sep=7pt, align=center, font=\footnotesize\sffamily, text=etInk
  },
  etBadge/.style={
    fill=#1, text=white, rounded corners=9pt, inner xsep=8pt, inner ysep=3pt,
    font=\scriptsize\bfseries\sffamily
  },
  etTitle/.style={font=\normalsize\bfseries\sffamily, text=etNavy, align=center},
  etSub/.style={font=\footnotesize\sffamily, text=etMuted, align=center},
  etArrow/.style={->, >=Stealth, line width=1.35pt},
  etDArrow/.style={<->, >=Stealth, line width=1.35pt}
}


\begin{document}
\begin{tikzpicture}[font=\sffamily]

% Title with >=10mm gap above panels
\node[etTitle] (title) at (0, 0) {Duas Lógicas na Criação de Novos Negócios (Sarasvathy, 2001)};

\def\panelW{7.2cm}
\def\panelH{7.65cm}
\def\stepW{5.9cm}

% ==================== LEFT PANEL: CAUSATION ====================
\node[draw=etNavy, fill=bgNavy, rounded corners=7pt, line width=1.25pt,
      minimum width=\panelW, minimum height=\panelH, anchor=north east]
  (leftPanel) at (-0.4, -1.25) {};

\node[etBadge=etNavy, anchor=north, inner xsep=10pt, inner ysep=4pt]
  (cBadge) at ([yshift=-10pt]leftPanel.north) {1. Lógica Causation (Planejamento Preditivo)};

\node[etBox=etNavy, text width=\stepW, minimum height=1.02cm, anchor=north]
  (c1) at ([yshift=-12pt]cBadge.south) {\textbf{Objetivo Definido}};

\node[etBox=etNavy, text width=\stepW, minimum height=0.88cm, anchor=north]
  (c2) at ([yshift=-14pt]c1.south) {\textbf{Previsão \& Análise}};

\node[etBox=etNavy, text width=\stepW, minimum height=0.88cm, anchor=north]
  (c3) at ([yshift=-14pt]c2.south) {\textbf{Alocação de Meios}};

\node[etBox=etNavy, text width=\stepW, minimum height=0.88cm, anchor=north]
  (c4) at ([yshift=-14pt]c3.south) {\textbf{Retorno Esperado}};

\draw[etArrow, etNavy] (c1.south) -- (c2.north);
\draw[etArrow, etNavy] (c2.south) -- (c3.north);
\draw[etArrow, etNavy] (c3.south) -- (c4.north);

\node[anchor=south, text width=6.4cm, align=center, font=\scriptsize\itshape\sffamily, text=etNavyDark]
  at ([yshift=10pt]leftPanel.south) {Ideal para mercados existentes e risco mensurável (plano clássico).};

% ==================== RIGHT PANEL: EFFECTUATION ====================
\node[draw=etGold, fill=bgGold, rounded corners=7pt, line width=1.25pt,
      minimum width=\panelW, minimum height=\panelH, anchor=north west]
  (rightPanel) at (0.4, -1.25) {};

\node[etBadge=etGold, anchor=north, inner xsep=10pt, inner ysep=4pt]
  (eBadge) at ([yshift=-10pt]rightPanel.north) {2. Lógica Effectuation (Incerteza Knightiana)};

\node[etBox=etGold, text width=\stepW, minimum height=1.02cm, anchor=north]
  (e1) at ([yshift=-12pt]eBadge.south) {\textbf{Meios Disponíveis}\\{\scriptsize\color{etMuted}(Quem sou · O que sei · Quem conheço)}};

\node[etBox=etGold, text width=\stepW, minimum height=0.88cm, anchor=north]
  (e2) at ([yshift=-14pt]e1.south) {\textbf{Perda Aceitável} {\scriptsize\color{etMuted}(não retorno máximo)}};

\node[etBox=etGold, text width=\stepW, minimum height=0.88cm, anchor=north]
  (e3) at ([yshift=-14pt]e2.south) {\textbf{Crazy Quilt} {\scriptsize\color{etMuted}(Parceiros comprometidos)}};

\node[etBox=etGold, text width=\stepW, minimum height=0.88cm, anchor=north]
  (e4) at ([yshift=-14pt]e3.south) {\textbf{Efeito Emergente}};

\draw[etArrow, etGoldDark] (e1.south) -- (e2.north);
\draw[etArrow, etGoldDark] (e2.south) -- (e3.north);
\draw[etArrow, etGoldDark] (e3.south) -- (e4.north);

\node[anchor=south, text width=6.4cm, align=center, font=\scriptsize\itshape\sffamily, text=etGoldDark]
  at ([yshift=10pt]rightPanel.south) {Ideal para deep tech nascente, spin-offs e ecossistemas incertos.};

\end{tikzpicture}
\end{document}
